\documentclass[11pt]{article}
\usepackage{icmlko}
\begin{document}

\title{압축된 라벨 --- 무리로 보면 보이는 것이 행으로는 안 보인다}
\author{월드모델 실험실}
\date{노트 484 · 2026년 8월 3일}
\maketitle

\begin{abstract}
노트 483 이 부정으로 끝났다 --- \emph{``행이 원래 어렵다''를 우리가 잰
어떤 변수로도 못 짚는다.} 그래서 모형을 놓고 \textbf{레코드를 열었다}.

\textbf{첫 줄에서 나왔다.} \texttt{venue\_prominence} 가 없는 $487$행은
\texttt{production} 이 \textbf{$100\%$ 비어 있고} 있는 $119$행은 채워져
있다. 그 축의 유무는 작품 성질이 아니라 \textbf{수집 사정} --- 어느
레코드에 풍부한 출처가 붙었나 --- 이다. 따라오는 것들도 그렇다: 태그
중앙 $8$ 대 $5$ · 갈래 $3$ 대 $2$ · 리뷰 수 $76$ 대 $61$.

그리고 하나가 더 있었다. \textbf{라벨 SD 가 $0.586$ 대 $0.721$ 로 좁다.}

\begin{center}
\begin{tabular}{lrrr}
\hline
무리 & $n$ & 라벨 SD & rho \\
\hline
\texttt{venue} 없는 $487$ & $487$ & $0.721$ & $+0.5877$ \\
$487$ 중 무작위 $440$ & $440$ & --- & $+0.5881$ \\
\textbf{$487$ 중 SD$\cdot$평균 맞춘 $440$} & $440$ & $\mathbf{0.570}$ & $\mathbf{+0.5265}$ \\
\texttt{venue} 있는 $119$ & $119$ & $0.586$ & $+0.4894$ \\
\hline
\end{tabular}
\end{center}

\textbf{라벨 분포만 맞추면 격차의 $62\%$ 가 사라진다.} 같은 크기 무작위는
$+0.5881$ 로 그대로이니 \emph{크기 탓이 아니다}.
\end{abstract}

\begin{figure}[t]
\centering
\includegraphics[width=\textwidth]{figs/f1.pdf}
\caption{\textbf{가} --- 무리의 라벨 SD 와 그 무리 안 rho. \textbf{나} ---
$-0.098$ 의 분해.}
\end{figure}

\section{노트 483 과 어긋나지 않는다}

노트 483 은 \emph{행 단위}로 극단성을 재서 $|$rho$|\ 0.088$ 을 얻고
``못 짚는다''고 적었다. 이 노트는 \emph{무리 단위}로 라벨 퍼짐을 맞춰
$62\%$ 를 설명한다. \textbf{둘은 같은 사실의 두 얼굴이다.}

행 하나하나는 자기 라벨이 가운데인지 끝인지로 조금밖에 안 갈리지만,
\textbf{무리 전체의 라벨이 좁으면 그 안에서 순위를 매길 거리가 통째로
줄어든다.} 앞엣것은 상관으로 재고 뒤엣것은 \emph{퍼짐}으로 재야 한다.

\begin{quote}
\textbf{조항} --- 무리끼리 rho 를 견줄 때는 \textbf{라벨 퍼짐을 맞추고
견준다.} 안 맞추면 \emph{잴 것이 적은 무리}를 \emph{모형이 못하는 무리}로
읽는다.
\end{quote}

\section{이것이 노트 470 에 걸린다}

노트 470 이 \textbf{무리 안 판}을 만들었다 --- 도메인을 갈래로 갈라
$0.4969 \to 0.4377$. 그때 대조는 \emph{같은 크기 무작위}였고($0.4932$),
그래서 ``크기가 아니라 동질성''이라 읽었다.

\textbf{그 동질성의 상당 부분이 라벨 압축일 수 있다.} 갈래 무리는 크기만
작은 게 아니라 \emph{라벨도 좁다} --- 무작위 쪼개기는 전체 퍼짐을 그대로
보존한다. 이 노트의 애니 사례에서 그 몫이 $62\%$ 였다.

\textbf{그러면 노트 470 의 읽기가 바뀐다.} ``판이 무리 사이 순위에
기댄다''(모형을 탓하는 읽기)에서 \textbf{``무리 안에서는 잴 것이
적다''}(자료를 적는 읽기)로. 무리 안 판이라는 자 자체는 남지만
\emph{그 낮은 값을 실력 부족으로 읽으면 안 된다.}

\textbf{안 쟀다.} 다음 노트의 일이다 --- 갈래 무리마다 라벨 퍼짐을
맞춘 대조를 놓고 $-0.0555$ 중 얼마가 압축인지 재야 한다.

\section*{남은 것}

\begin{itemize}
\item \textbf{노트 470 을 라벨 퍼짐 맞춰 다시} --- 위에 적은 그대로.
이 사이클에서 만든 자 하나의 뜻이 걸려 있다.
\item \textbf{남은 $38\%$} --- 라벨 퍼짐을 맞춰도 $119$행이 $+0.5265$
대 $+0.4894$ 로 낮다. 아직 모른다.
\item \textbf{\texttt{production} 이 붙은 레코드가 왜 라벨이 좁나} ---
많이 기록된 작품끼리는 서로 비슷하다는 뜻일 수 있다. 수집 갈래.
\end{itemize}

\end{document}
